% Transcription — fonds Alexandre Grothendieck, shelfmark 154, batch 2 (pages 21-40).
% Established from the facsimile archives/batches/154/batch-02.pdf, cut from a
% copy of 154.pdf fetched through the project's own relay
% (https://grothendieck-relay.onrender.com/source/154.pdf); PDF page k =
% archivists' page 20+k, no cover sheet. Per the skill transcribe-grothendieck.
% The folder is 175 pages in nine batches (1-20, ..., 161-175), transcribed in
% parallel; batch 1 was read first for notation and for the run it continues.
% Notation follows batch 1 and the transcription of folder 80.
%
% Pass: Opus 5.5 (claude-opus-5-5), 2026-09-26 — first pass, unchecked against the pages by a human.
%
% Dating and title. \dating is the folder's own date, copied verbatim from
% src/content/catalogue.ts; the folder belongs to the inventory group
% « Espaces stratifiés » (150-151), but since it has a date of its own the
% group's range is not added. Title from the catalogue, trailing full stop
% dropped; « [Système de pseudo-droites] » is bracketed there because the
% archivists supplied it.
%
% Paper. All twenty leaves are fan-fold computer-listing paper written in
% landscape, two columns to a leaf; the printouts (1982 job banners and
% accounting lines) show through or are written over. Not described further.
%
% Pages skipped: none. Pages summarised: none. Pages given only as a figure
% described in a French \note: 26 and 29 (coloured drawings, no text).
% Pages transcribed from his hand: 21-25, 27-28, 30-40 (figures described in
% notes).
%
% His pagination and dates. Pages 21-25 carry his numbers 10-14, continuing
% the run dated « 23.12.83 » that batch 1 records on pp. 11-19 (numbered
% 1-9; p. 20 a drawing without number). From p. 26 on no number of his is
% visible. Dates in his hand, inside that run: « 24.12. » at the head of an
% alinea on p. 25, « 25.12 » (underlined) on p. 27. Numbered/lettered items:
% a) b) c) (the first read as alpha) on p. 23; b) c) on p. 24; a) b) on
% p. 34; a) b) c) on p. 35; a) b) c) on p. 36; the four cases 1)-4) on p. 39
% and a) b) there. All recorded in \note{}s where they matter.
%
% What the batch is about.
%   21-25  (end of the 23.12.83 run) The polygon Pol(D) and its enlargement
%          Pol*(D), including pseudo-specialisation triangles (antipodes of
%          specialisation triangles) and free edges; weights of edges and
%          corners, possibly zero; the specialisation maps
%          Pol*(D) -> Pol*(D'), not isomorphisms in general, and the additive
%          relation of weights on p. 23; types of vertices of Pol*(D); the
%          observation that the weighted oriented cellular structure of
%          Gamma*(Sigma) does not distinguish pseudo-specialisations, but that
%          the specialisation maps are deducible (« Mais bien sûr ») via the
%          unfolded polygons, and that a 2-cell with weighted edges and
%          corners of fixed total weight gives a local system of unfolded
%          polygons on the underlying surface. 24.12: the dual map is the
%          good one; Pol(D) as the boundary of a small tubular neighbourhood,
%          the disc Delta(D), and the gluing of the discs Delta(D).
%   27-28  Gluing the Delta(D) without identifications gives a surface with
%          boundary, whose holes correspond to the vertices of the closed
%          surface. 25.12: Pol(Delta) for Delta in Sigma a supercritical
%          position; the two banks of Delta - {s}; an isomorphism from the
%          swept bank to I(D', s).
%   30-33  Abstract construction: a combinatorial graph with a
%          « semi-circularisation » (polygonal structures on the stars) and
%          transport isomorphisms phi_a between the orientation sets; the
%          surface obtained by gluing the dual discs Dis_s; faces as orbits of
%          rho_f; the operations sigma_0, sigma_1, sigma_2 on flags
%          (repères) and rho_som = sigma_2 sigma_1, rho_face = sigma_1
%          sigma_0, rho_f = sigma_2 sigma_0; the surface with boundary from
%          the doubled polygons Dis'_s, boundary components in bijection with
%          faces of C; for pseudo-line systems faces of order dividing 4, but
%          the privileged edge depends on the corner, so no splitting of C
%          into two dual graphs.
%   33-37  For a relative position L: the Möbius band B_L, the disc Delta_L,
%          the chord system (d_i) on it; axioms a) (two chords cross at most
%          once, inside) and b) (parallelism is an equivalence, classes of
%          nu >= 2 split into two antipodal segments), claimed necessary and
%          sufficient; arcs of generisation, of ordinary specialisation
%          (triangles adjacent to the boundary, conditions a)-c)), and
%          exceptional specialisation arcs (couloirs); two exceptional arcs
%          iff L is the « bissectrice » of a couloir; transition arcs, arcs
%          of inertia, transition vertices, three types of arcs of L~,
%          « rétréci » exceptional arcs; generisation as « écarter les
%          brins » (a décollement of the vertex s).
%   37-40  Attaching Delta_L to Delta_L' along corresponding generisation /
%          specialisation arcs, first by the evident identification and then,
%          on reflection, by the symmetry exchanging D and D'; all faces then
%          of order 4, coins with four sides of types spécialisation /
%          générisation / mixte (1)-4) on p. 39); faces in bijection with
%          pairs (L, a); what remains: the surface with boundary left by the
%          supercritical positions, the circuits of exceptional
%          specialisation arcs, and the start of the case L = D_i in Sigma.
%
% Notation fixed once for the batch, following batch 1: \underline{\Sigma}
% for his doubly underlined Sigma; \operatorname{Pol}(D),
% \operatorname{Pol}^{*}(D) (his superscript is a star or an alpha-like
% sign, unified as a star), \operatorname{Pol}^{\wedge}(D) on p. 25,
% \widehat{\operatorname{Pol}}; \Gamma(\Sigma), \Gamma^{*}(\Sigma);
% \widetilde{D}, \widetilde{L} for double covers / boundaries;
% \mathcal{C}, \mathcal{C}^{*} for the map and its dual; \operatorname{Dis}_s;
% \vec{a}, \vec{A}_s; \underline{\omega} for his underlined omega (sets of
% orientations); \Delta_L, B_L; « p.r. » position relative.
%
% Readings to check first: pp. 21-25 (fast, many \ill{}); p. 28 (a large
% cancelled box); the formulas of p. 31 (sigma_1, sigma_2 and the tilded
% operations); b) on p. 34; c) on p. 35; the end of p. 36; « antibalayage »
% and « décollement » on p. 37; « arc de neutralité » on pp. 37, 39.
\documentclass[11pt,a4paper]{article}
% Le préambule commun à toutes les transcriptions du fonds Grothendieck.
%
% Il tient en une idée : **l'incertitude se compose, elle ne se résout pas.**
% Une transcription qui lisse un mot douteux a détruit la seule chose pour
% laquelle on la faisait — un lecteur doit pouvoir savoir, à chaque endroit,
% ce qui était sur le feuillet et ce qui a été ajouté. Les six macros qui
% suivent sont donc l'appareil critique, et non de la décoration.
%
% Les mêmes commandes sont reconnues par `scripts/render.mjs`, qui produit la
% vue de lecture du site. Une macro ajoutée ici doit l'être aussi là-bas,
% faute de quoi le rendu échouera bruyamment — ce qui est voulu : un rendu
% silencieusement faux est pire qu'un rendu absent.

\ProvidesPackage{grothendieck}[2026/08/08 Transcriptions du fonds Grothendieck]

\usepackage{amsmath,amssymb,amsthm}
\usepackage{mathrsfs}
\usepackage{stmaryrd}
% Les diagrammes commutatifs sont partout dans ce fonds : c'est la langue
% même de la géométrie algébrique de Grothendieck, et une transcription qui
% les rendrait en prose ne transcrirait rien.
\usepackage{tikz-cd}
% Français par défaut — c'est la langue des feuillets — mais l'anglais est
% chargé aussi : la traduction et le résumé sont des documents anglais, et la
% typographie française (espaces avant les deux-points) y serait une faute.
% Ces documents basculent par \selectlanguage{english} en tête de corps.
\usepackage[english,french]{babel}
\usepackage{fontspec}
\usepackage{geometry}
\geometry{a4paper,margin=25mm,marginparwidth=28mm}
\usepackage{marginnote}
\usepackage{xcolor}
% ulem, not soul. `soul`'s \st and \ul are built on a character-by-character
% scan that XeLaTeX's font machinery breaks: the document fails with an
% undefined control sequence rather than degrading. ulem does the same two jobs
% and is Unicode-engine safe. `normalem` stops it hijacking \emph, which the
% transcriptions use for Grothendieck's own emphasis.
\usepackage[normalem]{ulem}
\usepackage{hyperref}
\hypersetup{colorlinks=true,linkcolor=black,urlcolor=black,pdfborder={0 0 0}}

% --- Métadonnées du lot -----------------------------------------------------
% Reprises telles quelles par le script de rendu, qui les affiche en tête de
% la vue de lecture. Elles fixent ce que le fichier prétend couvrir : sans
% elles, une transcription flotte sans point d'attache dans un fonds de seize
% mille feuillets.

\newcommand{\folder}[1]{\gdef\grothFolder{#1}}
\newcommand{\batch}[1]{\gdef\grothBatch{#1}}
\newcommand{\leaves}[2]{\gdef\grothFirst{#1}\gdef\grothLast{#2}}
\newcommand{\foldertitle}[1]{\gdef\grothFoldertitle{#1}}
\gdef\grothFolder{?}\gdef\grothBatch{?}\gdef\grothFirst{?}\gdef\grothLast{?}\gdef\grothFoldertitle{}

% --- Le résumé --------------------------------------------------------------
%
% Il ouvre la lecture modernisée, sous le titre « Résumé », et il est composé
% en petit. Ce n'est pas un ornement : c'est le seul passage du document qui
% ne soit pas des mathématiques, et un lecteur doit pouvoir le reconnaître
% avant de l'avoir lu — puis le sauter s'il connaît déjà le sujet, ou s'y
% arrêter s'il ne le connaît pas. Le filet qui le suit marque la reprise du
% corps.

\newenvironment{resume}
  {\small\setlength{\parskip}{0.45em}}
  {\par\medskip\hrule\medskip}

% --- Mots-clés ---------------------------------------------------------------
%
% En anglais, en clôture du résumé de la lecture modernisée : le vocabulaire
% moderne sous lequel chercher ce que les feuillets construisent. Le manifeste
% du site (`npm run manifest`) les extrait pour étiqueter la cote entière —
% cette ligne est la source unique des tags, il n'y a pas de fichier à côté.
\newcommand{\keywords}[1]{\par\smallskip\noindent
  {\footnotesize\textsc{Keywords} --- \textit{#1}}\par}

% --- L'appareil critique ----------------------------------------------------

% Le numéro de feuillet, en marge. C'est l'ancre qui permet de régler un
% désaccord sur une formule en regardant *un* feuillet plutôt qu'une chemise
% de six cents. Le site s'en sert aussi pour tourner les pages du fac-similé
% quand on fait défiler la transcription.
\newcommand{\leaf}[1]{%
  \marginnote{\footnotesize\color{gray!70}\textsf{#1}}%
  \label{leaf:#1}\ignorespaces}

% Illisible : on ne devine pas. Un mot inventé qui se lit comme les autres est
% le pire résultat possible.
\newcommand{\ill}{\textcolor{red!60!black}{[\,\ldots\,]}}

% Lecture douteuse : le mot est donné, et signalé comme incertain.
\newcommand{\uncertain}[1]{\uline{#1}}

% Ajout de l'éditeur — une lettre manquante, un mot sous-entendu.
\newcommand{\add}[1]{\textcolor{blue!50!black}{[#1]}}

% Biffé par Grothendieck lui-même. Ce qu'il a rayé fait partie du document :
% une rature dit souvent où la pensée a changé de direction.
\newcommand{\struck}[1]{\sout{#1}}

% Note du transcripteur, sur l'état matériel du feuillet.
\newcommand{\note}[1]{\par\smallskip{\small\color{gray!60!black}\textit{[#1]}}\par\smallskip}

% Note marginale de Grothendieck, distincte du corps du texte.
\newcommand{\marginal}[1]{\par\smallskip\noindent\hspace{1em}%
  {\small\color{gray!60!black}#1}\par\smallskip}

% --- Titre ------------------------------------------------------------------

\AtBeginDocument{%
  \begin{center}
    {\large\bfseries Fonds Alexandre Grothendieck --- cote n\textsuperscript{o}~\grothFolder}\\[2pt]
    {\itshape\grothFoldertitle}\\[4pt]
    {\small feuillets \grothFirst--\grothLast\ (lot \grothBatch)}\\[2pt]
    {\footnotesize\color{gray!60!black}Transcription établie d'après le fac-similé numérisé
     par l'Université de Montpellier.\\
     Les lectures incertaines sont soulignées, les illisibles notées [\,\ldots\,],
     les ajouts entre crochets.}
  \end{center}
  \vspace{1em}}

\endinput


\folder{154}
\batch{2}
\pages{21}{40}
\dating{1983-1984}
\watermark{Édition de démonstration}
\foldertitle{[Système de pseudo-droites] : notes manuscrites (1983-1984, s.d.)}

\begin{document}

\page{21}
\note{numéroté \emph{10} en haut à gauche : suite de la série commencée page 11 sous la date \emph{23.12.83} (numérotée 1 à 9 aux pages 11 à 19, la page 20 étant un dessin sans numéro)}
\note{figure, en haut à gauche : un arc (en vert) longé de part et d'autre par deux arcs bruns, traversé par six faisceaux de pseudo-droites orange ; sur les arcs bruns, des segments épais noirs marquent les arcs découpés à chaque croisement}

Les réflexions en termes des $P(D) \to P(D')$ et \ill{} un \add{(ou des triangles $T$)} \ill{} par spécialisation des triangles de \add{(s'il y en a)} \ill{}, \ill{} \ill{} \ill{} $T$ spécialisation (\ill{}, ou \ill{}) \struck{\uncertain{paraît}} \add{\uncertain{semble}} disparaître comme telle --- il faudrait donc indiquer la relation d'\emph{adjacence} des triangles de \struck{\ill{}} $D$), spécialisation de \ill{} comme \ill{} partie de $\Gamma(\Sigma)$ --- la structure \ill{}
\marginal{\emph{NB} elles \ill{} déduites \ill{} \ill{} structure de pondération}
\note{la note marginale est écrite en biais dans la marge de droite, séparée du texte par un trait}

\note{figure, à gauche : trois pseudo-droites orange sur une droite verte ; en pointillé vert, une générisation qui passe au-dessus du point de concours ; des traits noirs, pleins et pointillés, marquent les arcs découpés}
et la structure supplémentaire \ill{} \ill{} $\operatorname{Pol}(D)$ \ill{} en enlevant d'abord de $\operatorname{Pol}(D)$ les triangles de spécialisation (s'il y en a) \ill{} \ill{} : $T$, soit $\operatorname{Pol}(D;\,T)$ \struck{\ill{}} qui \ill{} \ill{} \uncertain{polygone}. On \ill{} \ill{} \ill{} \ill{} \struck{\ill{}} $\operatorname{Pol}(D)$, \struck{\ill{}} au contraire, augmenté \struck{$\operatorname{Pol}(D')$}, en y incluant des arcs des \uncertain{générisations} \ill{} des « pseudo-triangles de spécialisation » \ill{}
\marginal{\ill{} \ill{} \ill{} \ill{}}
\note{la note marginale, écrite en biais dans la marge de gauche, est reliée au texte par un trait vertical ; le « $\operatorname{Pol}(D;\,T)$ » est d'une lecture \uncertain{douteuse}}

\page{22}
\note{numéroté \emph{11} en haut}
un des \uncertain{extrémités} serait un sommet de $\Sigma$, comme
\note{figure : deux faisceaux de pseudo-droites orange sur une droite verte, un trait épais noir marquant sur celle-ci l'arc intercepté par chacun}

(ou bien \ill{} \ill{} \uncertain{couronne} d'une arête libre de $\Gamma(X)$, sans autre sommet au bout) ; \struck{il serait pratique} d'adjoindre, \add{tu} comme on dit plus haut, les antipodiques des \ill{} \ill{} (des \ill{} spécialisations, --- des pseudo-spécialisations --- pour \ill{} \ill{} des vraies spécialisations, \ill{} \struck{\ill{} \ill{} \ill{}} l'antipodique \ill{} \ill{} de spécialisation, \ill{} des pseudo-spécialisations) --- également des bords libres. Soit $\operatorname{Pol}^{*}(D)$ ce polygone (plus grand que celui précédemment désigné sous ce vocable). On trouve \ill{} \ill{} \ill{} \struck{\ill{}} \ill{} de $\operatorname{Pol}^{*}(D)$ \ill{} \ill{} \ill{} \ill{} \emph{\struck{nouveaux polygones}} de $\operatorname{Pol}^{*}(D')$, mais ce n'est pas \ill{} un isom., car des \ill{} \add{(ou triangles)} des pseudo-spécialisations peuvent être créés par spécialisation, dans le cas où \ill{} l'autre des deux faces adjacentes au triangle $T$ \ill{} l'une des \ill{} \ill{} \ill{} est un \uncertain{couloir}. Donc il faut \ill{} encore \ill{} \ill{} \ill{} \ill{} \ill{} antipodiques \ill{} \ill{} de l'\ill{} de
\[
\operatorname{Pol}^{*}(D) \to \operatorname{Pol}(D).
\]
\note{figures, colonne de droite : un faisceau orange sur une droite verte, les deux triangles de part et d'autre hachurés au crayon et marqués $t$, $t'$, le grand triangle $T$ au-dessous, une générisation en pointillé vert ; au-dessous, deux petits faisceaux orange sur une droite verte avec leurs arcs noirs ; plus bas, une petite figure analogue légendée $\operatorname{Pol}^{*}$}

\emph{NB} La pondération d'un \ill{} de $\operatorname{Pol}^{*}(D)$ peut être nulle --- \struck{\ill{}} entre \struck{deux \ill{} spécialisations} \add{\ill{}} générisations, \ill{} \ill{} \ill{} \ill{} \ill{} \ill{}.
\note{« $\operatorname{Pol}^{*}$ » : l'exposant est tantôt une étoile, tantôt un signe proche de $\alpha$ ; nous l'unifions en $^{*}$ sur la page, la graphie restant \uncertain{douteuse}. La phrase de la \emph{NB} se poursuit en une accolade dans la marge de droite (« \ill{} des ps.\ spéc.\ adjacents \ill{} figures \ill{} entre l'antipodique \ill{} arcs de spécialisation \ill{} de spécialisation »), qui se lit mal}

L'introduction de $\operatorname{Pol}^{*}(D)$ est \uncertain{intéressante}, \ill{} \ill{} \ill{} stable par antipodisme. \ill{} \ill{} \ill{} les compatibilités \ill{} \ill{} \ill{} \ill{} \ill{} \ill{} les pondérations, \ill{} dire que \ill{} \ill{} \struck{\ill{}} qui se correspondent \ill{} deux autres (\ill{}) pondérations, ainsi que \ill{} compatibles. Quand $\operatorname{Pol}^{*}(D) \to \operatorname{Pol}^{*}(D')$ \ill{} que se correspondent \ill{} n'est pas un iso, \ill{} il y a \struck{\ill{}} un \ill{} \ill{} deux des sommets adjacents \struck{\ill{}} \ill{} \ill{} \ill{} de $\operatorname{Pol}^{*}(D')$ qui \ill{} \ill{} $(D', a)$, $(D', a')$ \ill{} \ill{}.
\marginal{\emph{NB} \ill{} \ill{} \ill{} $\operatorname{Pol}^{*}(D)$ \ill{} \ill{} \ill{} \ill{}}
\note{note écrite en biais dans la marge de gauche}

\page{23}
\note{numéroté \emph{12} en haut. Figure en tête de la colonne de gauche : un arc portant, dans l'ordre, les points $x$, $b$, $a$}
\ill{} $b$ \ill{} \ill{} \ill{} $\operatorname{Pol}^{*}(D)$, \ill{} \ill{} l'arc $\widehat{ax}$ de $\operatorname{Pol}^{*}(D')$ ($x$ étant l'autre sommet de $\operatorname{Pol}^{*}(D')$ adjacent à $b$) on doit avoir
\[
\text{pondération de } \widehat{a,x} \text{ dans } \operatorname{Pol}^{*}(D') = \text{somme des pondérations de } \widehat{x,b},\ b,\ \widehat{b,a} \text{ dans } \operatorname{Pol}^{*}(D).
\]
\note{la formule est écrite sur trois lignes ; le premier $\operatorname{Pol}^{*}$ porte ici $(D)$, lu $(D')$ d'après la phrase précédente --- lecture \uncertain{probable}}

[Pour exprimer commodément cette relation, il faudrait développer la notion de pondération dans un sous-polygone d'un polygone pondéré\ldots, \ill{} telle façon que les \uncertain{déploiements} \add{pondérés} initiaux \ill{} \ill{} sous-polygone, soient \ill{} \ill{}.]

La question se pose si les \struck{\ill{}} \ill{} $\operatorname{Pol}^{*}(D) \to \operatorname{Pol}^{*}(D')$ peuvent se déduire des \ill{} structures de $\Gamma^{*}(\Sigma)$ (\struck{structure} cellulaire \ill{} et \struck{\ill{}} pondération).

Il faudrait d'abord pouvoir \ill{} les structures des polygones $\operatorname{Pol}^{*}(D)$ (\ill{} \uncertain{dehors} $D$ non supercritique).

a) Parmi ces sommets, \struck{\ill{} \ill{}} i.e.\ les arêtes de $\Gamma^{*}(\Sigma)$ en $D$, il y a les sommets de \emph{spécialisation} (arêtes \emph{\uncertain{partant}} de $D$), les arêtes de générisation (\ill{} \ill{} arrivant à $D$), et enfin \ill{} \ill{} \ill{} les \uncertain{arêtes} \ill{} \ill{} \uncertain{libres} en $D$), les sommets \ill{} \ill{} \ill{} \ill{} \ill{} \ill{} \ill{} \ill{} \ill{}. Notons \ill{} \ill{} « $b$ » \ill{} \ill{} \ill{} (\ill{}) \ill{} \ill{} :

$\alpha$) l'un des sommets \add{(\ill{} arêtes)} \struck{\ill{}} des générisations \ill{} \ill{} l'antipodisme ;

b) \ill{} \ill{} \ill{} \add{sommet} \struck{\ill{}} des spécialisations, \ill{} antipodique \ill{} \ill{} \add{sommet} \struck{\ill{}} des spécialisations ;

c) \ill{} \add{sommet} \struck{\ill{}} qui \ill{} un \struck{\ill{}} de générisation, \ill{} antipodique d'un \ill{} des spécialisations, \add{\ill{} \ill{}} \ill{} \ill{} \ill{} ps.\ spécialisation, \ill{} \ill{} dit antipodique d'un \ill{}, \ill{} \ill{} \ill{} qu'il est adjacent à \ill{} sommet de générisation, et la pondération de l'arête qui les joint est nulle ; \ill{} \ill{}
\note{la première rubrique, notée « a) » plus haut dans la phrase, est reprise par une lettre peu nette que nous lisons $\alpha$) ; les rubriques b) et c) sont écrites en marge de gauche de la colonne. Figure au bas de la page : deux pseudo-droites orange sur une droite verte, deux traits noirs marquant les arcs de part et d'autre du croisement}

\page{24}
\note{numéroté \emph{13} en haut}
Il ne semble pas que la structure de graphe orienté cellulaire pondéré de $\Gamma^{*}(\Sigma)$ permette de distinguer, parmi les \struck{sommets} arêtes issues d'un sommet $D$ i.e.\ les sommets de $\operatorname{Pol}^{*}(D)$, celles qui correspondent à des triangles de ps.\ spécialisation (parmi les \struck{paires \ill{} \ill{} à des arêtes de $\operatorname{Pol}^{*}(D)$ de pondération nulle, \ill{} antipodiques \ill{}} \add{sommets \ill{} incidents} \add{\ill{} \ill{}} \ill{} sont sommets de générisation).
\marginal{\ill{}, \ill{} !}
\note{la note marginale, écrite en biais entre les deux colonnes et séparée par un trait, se lit mal}

Cette distinction peut se déduire \ill{}, disposer des \ill{} de spécialisation \struck{\ill{} \ill{}} $\operatorname{Pol}^{*}(D) \to \operatorname{Pol}^{*}(D')$ : les sommets des ps.\ spécialisation sont distingués, parmi les « spécialisations » précédentes, comme les \struck{\ill{}} spécialisations \uncertain{\ill{}} des sommets de spécialisation, \uncertain{exactement}.

La prochaine question est donc si les \ill{} de spécialisation
\[
\operatorname{Pol}^{*}(D) \longrightarrow \operatorname{Pol}^{*}(D')
\]
sont déductibles des structures précédentes. \emph{Mais bien sûr}, via l'\struck{\ill{}} \ill{} correspondante des polygones \uncertain{déployés}, car il suffit \ill{} \ill{} \ill{} \ill{} \ill{} \ill{} \ill{} \ill{} (\ill{} \ill{}) \struck{\ill{}} \ill{} :

b) de savoir \struck{l'\ill{}} \add{l'\ill{}} \ill{} \ill{} l'une des deux orientations, ce qui est connu, et

c) de savoir l'image du sommet \add{$(D, a)$} particulier de $\widehat{\operatorname{Pol}}^{*}(D)$, qui est \ill{} comme l'antipodique, dans $\widehat{\operatorname{Pol}}(D')$, du sommet $(D', a)$ !
\note{la rubrique « a) » n'est pas visible : la phrase précédente, dont la fin se lit mal, en tient lieu. Dans c), le second $\widehat{\operatorname{Pol}}$ est écrit sans étoile}

En fait, chaque fois qu'il y a une 2-cellule \add{(avec bords isolés i.e.\ arêtes libres)} dont les \struck{\ill{}} (vrais) sommets sont d'ordre $2$, et dont les arêtes et coins sont pondérés [\ill{} pondération \ill{}] \struck{\ill{}} \ill{} des arêtes \ill{}, \ill{} pondération totale de chaque sommet \struck{\ill{}} (somme des pondérations des \add{arêtes} \struck{\ill{}} et coins incidents) étant \uncertain{fixe}, \ill{} trouve des \ill{} de transport entre les polygones locaux déployés, qui forment donc un syst.\ local sur la surface \uncertain{sous-jacente} à la carte !

\page{25}
\note{numéroté \emph{14} en haut}
Dans le cas actuel, il se trouve que \add{\ill{} \ill{}} \struck{\ill{} \ill{}} \ill{} \ill{} \ill{} \ill{}, \ill{} \ill{} $\operatorname{Pol}^{*}(D) \to \operatorname{Pol}^{*}(D')$ \ill{} les propriétés \ill{}.

Il faudrait examiner à part la relation de spécialisation vers un $\Delta$ \uncertain{supersingulier}.

\emph{24.12.} Finalement, il apparaît que c'est le polygone $\operatorname{Pol}^{\wedge}(D)$ qui est le « bon », ou encore \ill{}, le polygone \ill{} \ill{} \ill{} \struck{\ill{}} au \ill{} \ill{} \ill{} \ill{} à la carte \emph{duale} de celle considérée \ill{} --- c'est la carte duale qui est la plus \uncertain{intéressante} \ill{} \ill{} \ill{} \ill{} \ill{} \ill{} \ill{} le polygone \struck{\ill{}} $\operatorname{Pol}(D)$ (\ill{} \ill{} !) comme \struck{le bord} associé au polytope, bord d'un petit voisinage tubulaire \ill{} de $D$, ou encore bord du disque $\Delta(D)$ \ill{} compte de \ill{} \ill{} voisinage tubulaire, \ill{} la carte qui \ill{} intérieur peut être considérée comme obtenue en recollant les disques $\Delta(D)$, d'une façon naturelle que je viens de m'expliquer \ill{} \ill{}, \struck{\ill{}} qui \ill{} déduits des \add{\ill{} disques $\Delta(D)$} ces disques en prenant pour bord les « arêtes libres ».
\note{« 24.12. » est écrit en tête de l'alinéa, en marge : date de sa main (24 décembre 1983, dans la suite datée \emph{23.12.83} en page 11). « $\operatorname{Pol}^{\wedge}$ » rend un exposant en forme de chevron}

\ill{} \ill{} que ces \ill{} \ill{} \add{(\ill{} \ill{} et \ill{} \ill{} des générisations)} spécialisations \ill{} et on les \ill{} \ill{} la relation, de sorte qu'on trouve un disque $\Delta'(D)$ dans lequel les \ill{} \add{\uncertain{sommets}} précédentes apparaissent comme segments \ill{} du bord vers l'intérieur (si on garde la trace du \struck{\ill{}} syst.\ de droites (\ill{} \ill{}) à l'intérieur de $\Delta'(D)$ (cf.\ figure ci-contre, voisinage du bord) au \ill{} des points doubles
\note{figures au bas de la page : un grand disque (bord vert) traversé par un arrangement de pseudo-droites orange, des régions coloriées en jaune, les points du bord numérotés $1$ à $12$ au crayon, de petits segments bleus marquant les extrémités ; à côté, un petit disque au bord vert et bleu, d'où partent vers l'intérieur dix « éventails » de traits orange et verts numérotés $1$ à $10$ --- le disque $\Delta'(D)$, vu au voisinage du bord}

\page{26}
\note{page de dessins, sans texte : des faisceaux de pseudo-droites orange coupant une droite verte (pleine ou en pointillé), avec leurs arcs noirs, deux marqués $b$ ; un point de concours entouré des étiquettes $b$, $b-1$, $b+1$, $b\pm 2$ \uncertain{lues} en partie à l'envers ; une ellipse verte ; au crayon, deux fuseaux traversés de droites ; une générisation en pointillé vert passant sous deux faisceaux ; enfin un point double traversé d'une flèche orange en pointillé (une spécialisation traversant le sommet)}

\page{27}
\ill{} \uncertain{artificiels} aux bords des disques, \ill{} \ill{} --- i.e.\ certaines ps.\ droites des $\Delta(D)$, qui ne se rencontrent pas dans $\Delta(D)$, vont se rencontrer sur le bord de $\Delta'(D)$ \struck{\ill{}}. Si on fait par contre les recollements \add{des} \struck{\ill{}} $\Delta(D)$, sans faire les identifications qui mènent aux $\Delta'(D)$, on trouve une \emph{surface à bord} (\add{bord} \ill{} \ill{}), d'où à nouveau, en bouchant \add{par des disques $\Delta_x$} les trous, une surface sans bord compacte. (On obtient la surface envisagée au début en contractant \add{\ill{} des} ces disques en \struck{\ill{}} un point.) La surface \uncertain{tubulaire} à bord obtenue est telle que les pts du bord sont \add{tous d'indice $2$ ou trois} \struck{doubles ou triples}, et chaque sommet est sur le bord, elle est \ill{} de la carte \ill{} bord par le procédé canonique, consistant à entourer un petit disque autour de chaque sommet. En tous cas, les trous à boucher dans la surface à bord correspondent canoniquement aux sommets de la surface sans bord.
\marginal{\ill{} \ill{} \ill{}}
\note{le numéro de page n'est pas visible en haut ; la note marginale, écrite en biais, renvoie par une flèche au début de la colonne de droite}

\emph{25.12}\quad Il faudrait regarder la relation entre $\operatorname{Pol}(D)$ et $\operatorname{Pol}(D')$, quand $D \to D'$ et $D'$ supercritique. Il faudrait déjà \emph{définir} $\operatorname{Pol}(D')$, i.e.\ $\operatorname{Pol}(\Delta)$ quand $\Delta \in \Sigma$, comme l'ens.\ des générisations de $\Delta$, avec une structure polygonale convenable dessus. Or on a vu que les \struck{\ill{}} générisations de $\Delta$ correspondent aux sommets sur $\widetilde{\Delta}$, image inverse de $\Delta$ dans $\widetilde{X}$. D'ailleurs pour un sommet $s$ de $\Delta$,
\note{« 25.12 », souligné, est une date de sa main. Figures : deux disques au bord vert, traversés de pseudo-droites orange ; le premier hachuré au crayon sur sa moitié gauche, le second avec deux régions triangulaires coloriées en jaune au bord}

\page{28}
les pts de $\widetilde{\Delta}$ \uncertain{au-dessus} de $s$ i.e.\ les orientations de $X$ en $s$, correspondent aux \add{deux} « rives » de $\Delta \smallsetminus \lbrace s \rbrace$ dans $X$. La spécialisation $D \to D'$ définit une de ces rives (celle balayée par la spécialisation), et on trouve un isom.\ combinatoire canonique entre le segment combinatoire correspondant à la rive des \struck{spécialisations} \add{\uncertain{générisations}} correspondant à la rive balayée \add{\ill{}} de $D$, et celui \ill{} i.e.\ le \emph{\uncertain{découpage}} de $D'$ \ill{} \ill{} $D'$ le \uncertain{droite} d'une $D' \smallsetminus \lbrace s \rbrace$, (En termes de $\widetilde{D}'$, \ill{} \ill{} un antipodique $\widetilde{s}^{-}$, \struck{point} \add{\uncertain{sommet}} $\widetilde{s}$ de $\widetilde{D}'$, d'où un \ill{} \add{deux orientations de} $\widetilde{D}'$ en deux \uncertain{arcs}, correspondant aux \struck{\ill{}} deux arcs de $\widetilde{D}'$ \uncertain{bordés} par \ill{} extrémités $\widetilde{s}^{-}$). \struck{\ill{} antipodisme \ill{} $\widetilde{D}'$ qui \ill{} \ill{} \ill{} \ill{} $\widetilde{s}$, $\widetilde{s}^{-}$ \ill{} \ill{} \ill{}} orientés d'origine $\widetilde{s}$, extrémité \ill{} \ill{} \struck{\ill{}} combinatoire \struck{\ill{} \ill{} \ill{}} \ill{} \ill{} deux arcs, et qui \ill{} \ill{} \ill{} \ill{} \ill{} \ill{} \ill{}, \ill{} \ill{} \ill{} \ill{} \ill{} \ill{} \ill{} \ill{} combinatoire canonique\ill{} l'antipodisme définit un isom.\ canonique entre les deux arcs, qui, \ill{} \ill{} \ill{} combinatoirement, sont $\mathrm{I}(D', s)$ \ill{} \ill{} \ill{} \ill{} un segment combinatoire. \ill{} \ill{} \ill{} un \uncertain{isom.}
\[
\begin{array}{c} \text{rive des spécialisations} \\ \text{de } D \text{ \uncertain{utilisée}} \end{array} \xrightarrow{\ \sim\ } \mathrm{I}(D', s)
\]
\note{la moitié inférieure de la colonne est couverte d'un grand encadré barré de plusieurs traits obliques et de lignes biffées, avec des ajouts interlinéaires ; nous n'en donnons que ce qui se lit. « $\mathrm{I}(D', s)$ » rend un I majuscule appuyé ; la formule finale est écrite seule, en haut de la colonne de droite, le reste de la feuille étant blanc}

\page{29}
\note{page de dessins, sans texte : cinq disques au bord vert, chacun traversé d'un arrangement de pseudo-droites orange, accolés les uns aux autres en chaîne autour d'un point où les bords se rejoignent ; sur l'un des bords, un arc vert foncé fléché porte les lettres $p$, $q$ et un troisième signe \uncertain{peu net} ; hachures et gribouillis au crayon sur deux disques --- recollement des disques $\Delta(D)$ décrit pages 25 et 27}

\page{30}
Considérons un graphe comb.\ $\Gamma = (S, \vec{A}, \sigma_0, \alpha)$ avec une \emph{semi-circularisation} (structure polygonale \add{donnée} sur l'un des sommets issus d'un \uncertain{sommet} donné) et \struck{pour toute} d'isom.\ de transport
\[
\varphi_{\vec{a}} : \underbrace{\omega(\vec{A}_s)}_{\omega_s} \longrightarrow \underbrace{\omega(\vec{A}_t)}_{\omega_t} \qquad \text{pour } \vec{a} : s \to t
\]
(satisfaisant la seule condition
\[
\varphi_{-\vec{a}} = \varphi_{\vec{a}}^{-1}\ ).
\]
\note{« semi-circularisation » est souligné ; $\omega(\vec{A}_s)$ désigne, comme plus haut, l'ensemble des deux orientations de la structure polygonale de $\vec{A}_s$}

On suppose pour simplifier que les ordres des sommets sont $\geq 3$ (sinon il faut se donner, pour les sommets d'ordre $1 \leq \nu \leq 2$, une structure de polygone combinatoire $\supset \vec{A}_s$ \ill{} un de ses sommets\ldots, cont.\ un \ill{} des deux orientations, de façon à \ill{} \ill{} : la donnée des $\varphi_{\vec{a}}$). La surface associée aux $\varphi_{\vec{a}}$ \uncertain{s'obtient} ainsi : on considère la réalisation \add{par disque polygonal} topologique canonique du polygone dual du disque $\vec{A}_s$ (ce qui \ill{} \ill{} via une décomposition barycentrique de $\vec{A}_s$, qui donne une réalisation de $\overline{A}_s$, et pour la réalisation \struck{\ill{}} $\operatorname{Dis}_s$ \uncertain{du} polygone dual, chaque \uncertain{branche} comme réunion de $2\nu_s$ \emph{drapeaux} correspondant aux repères de $\vec{A}_s$ --- ici les
\note{« $\overline{A}_s$ » : la barre est ici horizontale, non fléchée ; la parenthèse ouverte n'est pas refermée sur la page}

\note{figure, en haut à droite : une carte à faces quadrangulaires tracée en vert foncé, son dual en vert clair, les faces hachurées au crayon ; un sommet $s$ marqué d'un point, une arête fléchée $\vec{a}$ de $s$ vers $t$, et les orientations $\omega_s$, $\omega_t$ figurées par de petites flèches courbes ; un cercle brun entoure un sommet de la carte}
\marginal{cas $\nu_s = 1, 2$ n'est pas exclu}
\note{en marge, en tête de la colonne de droite}
arêtes de ces polygones correspondant aux \struck{faces} \add{sommets} de $\vec{A}_s$ i.e.\ \struck{dans \uncertain{pol}} aux arêtes de $\Gamma$ d'origine $s$. \uncertain{Ceci} dit, on recolle, pour $\vec{a} \in \vec{A}$, \struck{$\vec{A}_s$} pour $s = \operatorname{or}(\vec{a})$, $t = \operatorname{ex}(\vec{a}) = \operatorname{or}(-\vec{a})$) \struck{l'arête de $\operatorname{Dis}_s$} l'arête $\alpha_{\vec{a}}$ de $\operatorname{Dis}_s$ correspondant à $\vec{a}$, avec l'arête $\alpha_{-\vec{a}}$ de $\operatorname{Dis}_t$ correspondant à $-\vec{a}$, moyennant un isom.\ \struck{\ill{}} \add{\ill{}} sur les arêtes (\uncertain{donc} défini \struck{\ill{}} \add{\ill{}} \struck{\ill{} \ill{}} \add{\ill{} de l'une} des deux est \uncertain{déterminée}, d'extrémités) \struck{\ill{} \ill{} \ill{}} à l'aide de l'isom.\ de transport $\varphi_{\vec{a}}$ : la \uncertain{règle}, c'est que si l'orientation $\omega_s$ de $\vec{A}_s$ i.e.\ de $\operatorname{Dis}_s$ et $\omega_t$ de $\vec{A}_t$ i.e.\ de $\operatorname{Dis}_t$ se correspondent, alors l'isom.\ d'identification
\[
\psi_{\vec{a}} : \alpha_{\vec{a}} \xrightarrow{\ \sim\ } \alpha_{-\vec{a}}
\]
doit \struck{\ill{}} \emph{renverser} les orientations induites par $\omega_s$ sur $\alpha_{\vec{a}}$ resp.\ $\omega_t$ sur $\alpha_{-\vec{a}}$.
\note{le nom de l'isomorphisme d'identification est lu $\psi_{\vec{a}}$, lecture \uncertain{probable}}

\page{31}
Par ces identifications, on trouve une surface topologique, \add{\uncertain{compacte} sans bord}, dans laquelle la réalisation topologique $|\Gamma|$ de $\Gamma$ est canoniquement plongée. Les faces orientées de cette surface se réalisent comme les orbites sous l'opération $\rho_f$ sur l'ens.\ \struck{\ill{}} $\operatorname{Rep}$ des couples $(\vec{a}, \omega)$, où
\[
\vec{a} \in \vec{A}, \quad \omega \in \underline{\omega}_s \overset{\mathrm{def}}{=} \underline{\omega}(\vec{A}_s),
\]
opération définie par \struck{$\rho_f(\vec{a}, \omega) =$}
\[
\left\lbrace
\begin{array}{l}
\rho_f(\vec{a}, \omega_s) = (\vec{b}, \omega_t) \\
\vec{b} = \rho_{\omega_t}^{-1}(-\vec{a})
\end{array}
\right.
\qquad
\begin{array}{l}
\text{où } \omega_t = \varphi_{\vec{a}}(\omega_s) \in \underline{\omega}(A_t) \\
\text{et où } b \in \vec{A}_t \ (t = \operatorname{ex}(\vec{a})) \\
\text{est donné par}
\end{array}
\]
\note{les accolades de sa main sont rendues par un tableau ; dans la première ligne, « $\omega_t$ » est écrit au-dessus d'un terme raturé}
où $\rho_{\omega_t}$ est la permutation circulaire sur $\vec{A}_t$ définie par l'orientation $\omega_t$. C'est aussi l'opération de rotation autour des sommets de la carte $\mathcal{C}^{*}$, définie par la façon en question.
\note{« $\mathcal{C}^{*}$ » : la carte duale, lecture \uncertain{probable} de son C majuscule suivi d'une étoile}

Si on veut une description combinatoire complète de la carte $\mathcal{C}$ définie par
\[
(\Gamma, (\pi_s)_{s \in S}, (\varphi_{\vec{a}})_{\vec{a} \in \vec{A}})
\]
il faut \struck{\ill{}} expliciter les opérations \add{\uncertain{fondamentales} sur $\vec{A}_s$} $\sigma_0, \sigma_1, \sigma_2$ sur $\operatorname{Rep} \mathcal{C}$, ens.\ des repères. On trouve
\[
\left\lbrace
\begin{array}{l}
\sigma_0(\vec{a}, \omega_s) = (-\vec{a}, -\varphi_{\vec{a}}\,\omega_s) \\
\sigma_2(\vec{a}, \omega_s) = (\vec{a}, -\omega_s) \\
\sigma_1(\vec{a}, \omega_s) = (\rho_{\omega_s}(a), -\omega_s)
\end{array}
\right.
\]
\note{$(\pi_s)_{s \in S}$ : la lettre est \uncertain{peu nette} (les structures polygonales données). Figure à gauche des formules : deux arêtes fléchées (vert foncé et vert clair) se croisant, $s$ au bas, $t$ en haut ; les quatre repères voisins, hachurés au crayon, sont marqués $r$, $\sigma_0 r$, $\sigma_1 r$, $\sigma_2 r$, une petite flèche courbe tournant autour de $s$}

Soit encore, en termes des repères $r$ des $\vec{A}_s$, en interprétant
\[
\operatorname{Rep}(\mathcal{C}) \simeq \coprod_{s \in S} \operatorname{Rep}(\vec{A}_s)
\]
\note{sous $\operatorname{Rep}(\vec{A}_s)$, accolade : « ens.\ des drapeaux », suivi d'un « \struck{ens.\ de sommets} $\vec{a}$ » biffé}
\[
\left\lbrace
\begin{array}{l}
\sigma_0(r) = r', \text{ défini par } \left\lbrace \begin{array}{l} \vec{a}_{r'} = -\vec{a} \\ \omega_{r'} = -\varphi_{\vec{a}}(\omega_r) \end{array} \right. \\
\sigma_1(r) = \widetilde{\sigma}_0(r) \\
\sigma_2(r) = \widetilde{\sigma}_1(r)
\end{array}
\right.
\]
où $\widetilde{\sigma}_0^{s}, \widetilde{\sigma}_1^{s}$ sont les opérations habituelles sur $\operatorname{Rep}(\vec{A}_s)$, $\widetilde{\sigma}_i^{s}$ changeant la composante \add{\uncertain{du repère}} de dimension $i$, pas l'autre.
\note{dans la première ligne, un premier second membre est biffé et récrit « $r'$ »}

En termes de ces opérations, on trouve
\[
\left\lbrace
\begin{array}{l}
\rho_{\mathrm{som}} = \sigma_2\sigma_1 \ \bigl(= \widetilde{\sigma}_0^{s}\widetilde{\sigma}_1^{s} = \widetilde{\rho}^{s}, \text{ opération de rotation sur les repères d'un polygone comb.}\bigr) \\
\rho_{\mathrm{face}} = \sigma_1\sigma_0
\end{array}
\right.
\]
\note{au-dessus de la parenthèse, en ajout : « sur $\operatorname{Rep}_s = \operatorname{Rep}(\vec{A}_s)$ »}
d'où
\[
\rho_f = \sigma_2\sigma_0 \overset{\mathrm{def}}{=} \sigma, \qquad \sigma(\vec{a}, \omega_s) = (-\vec{a}, \varphi_{\vec{a}}\,\omega_s).
\]
\note{« $\rho_f$ » est écrit après un premier « $\rho_{\mathrm{face}}$ » biffé ; la dernière formule est séparée par un trait vertical}

\page{32}
La subdivision barycentrique de $\mathcal{C}$, qui est aussi celle de $\mathcal{C}^{*}$, est apparente sur les dessins donnés, ses faces sont les drapeaux qui ont servi de composants pour la construction des $\operatorname{Dis}_s = |\vec{A}_s|$. (Cf traits \ill{} \ill{} \ill{})
\note{il renvoie à trois traits de couleur tracés dans la ligne : un trait vert foncé, un trait vert clair, un trait pointillé noir --- les arêtes de $\mathcal{C}$, de $\mathcal{C}^{*}$ et de la subdivision}

On va donner une autre construction, de la surface $\mathcal{X}$ sous-jacente à $\mathcal{C}$, qui s'obtient intuitivement en traçant des petits cercles \struck{\ill{}} autour des sommets de \struck{la face} \add{\uncertain{carte}} \add{(centres des faces de $\mathcal{C}$)} le \uncertain{dual}, et enlevant les portions d'arêtes de $\mathcal{C}^{*}$ (bleues) qui se trouvent à l'intérieur. La surface \emph{à bord} déduite de $\mathcal{C}^{*}$ en enlevant les disques, avec la décomposition cellulaire définie par la construction précédente (\ill{} les faces de $\mathcal{C}^{*}$ étaient réalisées par les $\operatorname{Dis}_s$ ($s \in S$)), s'obtient en remplaçant les $\operatorname{Dis}_s$ par les $\operatorname{Dis}'_s$, obtenus en « enlevant les coins » de ces plaques polygonales. Chaque
\note{« bleues » : les arêtes de $\mathcal{C}^{*}$ sont pourtant tracées en vert clair sur les dessins de cette page et de la page 30}

$\vec{A}_s$ est remplacé par le polygone combinatoire « doublé » $\vec{A}'_s$ \struck{(l'ens.\ des $2$ sommets)} \add{(dont les $2$ \emph{arêtes}} sont les repères de $\vec{A}_s$), \struck{dont le} d'ordre $2\nu$, et on prend le polygone doublé, dont les « arêtes » \add{(« d'insertion »)} sont en corr.\ biunivoque avec les \struck{repères} \add{arêtes de $\vec{A}_s$} et les « arêtes de recollement » avec les \add{sommets} de $\vec{A}_s$, \struck{pour tout tel repère, c'est les sommets} \add{\ill{}} de $\vec{A}'_s$. Ceci dit, on recolle les $\operatorname{Dis}_s$ ($s \in S$) suivant les arêtes de recollement seulement, on trouve une carte à bord dont les faces sont les $\operatorname{Dis}'_s$, et les comp.\ connexes du bord sont en corr.\ 1-1 avec les \struck{\ill{}} faces de $\mathcal{C}$ i.e.\ les sommets de $\mathcal{C}^{*}$ --- donc ce sont les orbites sous le groupe diédral $\langle \sigma_1, \sigma_0 \rangle$ contenant $\langle \rho_f \rangle$ comme sous-groupe d'indice $(1 \text{ ou})\,2$. Notons que dans le procédé de recollement précédent, chaque composante connexe \struck{d'ordre} du bord apparaît comme un polygone topologique, l'ordre \uncertain{égal} à l'ordre
\note{la seconde colonne commence en haut à droite de la feuille. Figures dans la marge de cette colonne : un disque au bord vert foncé, rayonné de quatre traits vert clair, avec de petits arcs bruns aux coins du bord (les coins enlevés) ; au-dessus, un petit polygone colorié en vert, marqué $s$ et $t$}

\page{33}
du sommet correspondant de $\mathcal{C}^{*}$. De façon plus précise, les arêtes correspondent biunivoquement aux \emph{arêtes} du polygone local des arêtes \struck{\ill{}} issues de $\sigma$, ou encore aux sommets de $\Gamma$ incidents à la face. Plus précisément encore, le polygone comb.\ associé à cette \struck{\ill{}} composante du bord est can.\ isom.\ au \emph{dual} du polygone combinatoire correspondant à la face envisagée de $\mathcal{C}$.

Dans le cas qui nous intéresse, provenant d'un syst.\ de ps.\ droites $\Sigma$, en prenant la carte \uncertain{d'}\ill{} des \struck{\ill{}} toutes ses positions relatives, il se trouve que l'ordre des \uncertain{faces} \add{\ill{}} ces faces est (un diviseur de) $4$. De plus, parmi les \add{quatre} arêtes de chaque comp.\ connexe du bord, correspondant donc à un \emph{coin} de la carte $\mathcal{C}$, il y en a une qui est privilégiée.
\note{le $\sigma$ de « issues de $\sigma$ » désigne vraisemblablement le sommet de $\mathcal{C}^{*}$ en question}

Mais la propriété de « type » qui la distingue ne dépend pas \add{(en général)} que du sommet correspondant de $\mathcal{C}$, \add{i.e.\ de $\Gamma$}, mais bien du coin \struck{qui le} particulier envisagé. Donc il n'y a pas question d'espérer, en général, de réaliser $\mathcal{C}$ comme \add{\uncertain{le graphe} \uncertain{composé}} réunion de deux \struck{cartes $\mathcal{C}_0$, $\mathcal{C}$} graphes $\gamma, \gamma^{*}$, correspondant respectivement à deux \struck{graphes de} \add{cartes} \uncertain{duales} l'une de l'autre, sur $\mathcal{X}$, \ill{} les \ill{} de \ill{} \ill{}.
\note{la fin de cet alinéa, écrite sur des lignes du listing, se lit mal}

Soit $\underline{\Sigma} = (X, \Sigma)$ un syst.\ de pseudodroites, et $L$ une position relative (non triviale). On prend une \struck{petite} \add{étroite} bande \add{(de Möbius) $B_L$} autour de $L$, dont le bord $\widetilde{L}$ peut être vu comme revêt.\ à deux feuillets de $L$ (\emph{cov.\ des rives}), et le complémentaire de $\mathring{B}_L$ est un disque, \struck{\ill{}} soit $\Delta_L$. Toute $D_i \in \Sigma$ coupe $\partial \Delta_L = \widetilde{L}$ en exactement deux points, et sa trace \struck{\ill{}} sur $\Delta_L$ est un segment (ainsi que sa trace sur $B_L$).
\note{« cov.\ » : lecture \uncertain{douteuse} (revêtement des rives) ; la lettre $L$ est écrite sur un $D$ biffé à sa première occurrence}

\page{34}
\note{figure en tête de la colonne de gauche : un disque au bord vert traversé de cordes orange ; sur le bord, six segments épais vert foncé, et six petits triangles coloriés en orange au voisinage du bord}
On trouve un « syst.\ de cordes » \add{$(d_i) = (D_i \cap \Delta)$} dans $\Delta = \Delta_L$ qui a les propriétés suivantes

a) Deux $d_i$ se coupent en au plus un point, lequel est \emph{intérieur} à $\Delta$, et elles se croisent en ce point.

(cela fait partie de la notion $=$ de « syst.\ de cordes », avec la condition que les $d_i$ soient des segments topologiques et $\partial d_i = d_i \cap \partial\Delta$)
\note{« admis » est ajouté au-dessus de « cordes »}

b) La relation dans $I$ « $d_i = d_j$ ou \struck{\ill{}} $\partial d_i \cap \partial d_j = \emptyset$ » (« parallélisme ») est une relation d'équivalence, et \struck{\ill{}} si on a une classe d'équivalence \struck{\ill{}} \add{\ill{}} $\nu \geq 2$ \struck{\ill{} \ill{} \ill{}}, alors l'ensemble des \add{2$\nu$} sommets correspondants est réparti en $2$ \add{segments} \add{$S', S''$} de $\nu$, [naturellement] \uncertain{antipodiques} [\emph{NB} ceci a un sens \uncertain{précis} combinatoire, les \uncertain{sommets} des syst.\ étant \uncertain{munis} \uncertain{rangés} ($\partial\Delta$, \ill{}) \ill{} \ill{} \ill{} $S \cap S'$ et $S \cap S''$ réduits à un point \ill{} $\partial\Delta$ ($\sim$), et les sommets de $\partial\Delta$ \uncertain{circulaires} formant un c.c.\ et de C \ill{} $(s_1, s_2, \ldots, s_\nu)$ i.e.\ de sommets consécutifs adjacents. [Cette partition de l'ens.\ des $2\nu$ sommets est \uncertain{unique}]
\note{b) est d'une écriture très rapide, sur les lignes du listing, avec des ratures ; les crochets sont de sa main et ne se referment pas tous}

Ces conditions combinatoires nécessaires pour que le syst.\ de cordes soit isomorphe à celui d'une position relative \uncertain{pour} \uncertain{un} $\Sigma$ convenable, sont aussi suffisantes. Elles impliquent que si on choisit un représentant dans chaque classe de parallélisme, on trouve un syst.\ de cordes strict (i.e.\ \uncertain{totalement} parallèle).
\note{« totalement parallèle » : la lecture de l'adverbe est \uncertain{douteuse}}

Pour $\mathcal{C}$ donné, considérons les arcs topologiques interceptés par les ensembles $S'$, $S''$ on trouve deux arcs (topologiques, ou combinatoires), appelés \emph{arcs de générisation} relatifs à $L$ (et à $\mathcal{C}$). L'ens.\ de tous ces arcs, pour $L$ variable, \add{\uncertain{divisé} par l'antipodisme,} (est en corr.\ biunivoque avec l'ensemble $L \cap S$ ($S = S(\Sigma)$).

On définit de plus les \emph{arcs de spécialisation}, en considérant les triangles formés par le système de cordes, \struck{adjacents} adjacents à $\partial\Delta$. On considère les ensembles de sommets
\note{le $\mathcal{C}$ de « Pour $\mathcal{C}$ donné » désigne vraisemblablement le système de cordes ; « arcs » et « générisation » sont soulignés}

\page{35}
$\sigma$ ayant les propriétés suivantes

a) \struck{il existe} \add{$\operatorname{card}\sigma \geq 2$, et} $\sigma$ forme un arc combinatoire \uncertain{dans} l'un des sommets du syst.\ de cordes.

b) Les $D_{i(s)}$ passant par les $s \in \sigma$ sont mutuellement disjoints, \struck{\ill{}} passant par un même sommet (\uncertain{bien} \uncertain{sûr} unique) \struck{\ill{}}, et toutes \struck{droites} \add{cordes} passant par $t$ sont parmi les $D_{i(s)}$.

c) \struck{Toutes} Les $s \in \sigma$ (\ill{} qui \uncertain{vont} $\rightrightarrows$, \struck{\ill{} \ill{} \ill{} \ill{}} les $s$ extrêmes) \uncertain{n'appartiennent} pas à \uncertain{un} arc de \struck{spécialisation} générisation, \ill{} que \uncertain{revient} au $\rightrightarrows$, ne sont pas \add{\uncertain{déjà}} \struck{\ill{}} un sommet $S$ de $\Sigma$.
\note{la troisième ligne de b) commence par un mot biffé ; c) porte une ligne entière biffée, que nous ne lisons pas. Figure au bas de la colonne : quatre pseudo-droites orange issues d'un point $t$ et coupant une droite vert clair aux points $s_0, s_1, s_2, s_3$}

Sous ces conditions, et supposant \uncertain{bien} \uncertain{fixé} \add{\uncertain{\ill{}}} \uncertain{(\ill{})} le système $\Sigma$ non trivial, \struck{alors} il existe pour \ill{} $s \in \sigma$ un unique couple d'extrémités $(t, s)$, et pour $s, s' \in \sigma$ adjacents, un unique triangle de \add{$\Sigma$ de} sommets $t, s, s'$. La réunion de tous ces triangles s'appelle un triangle de spécialisation. Sa base, qui est l'arc \uncertain{géom.}\ intercepté par l'arc combinatoire $\sigma$, est appelée un arc de spécialisation (ordinaire) \struck{\ill{}}. Tous ces arcs (\uncertain{générés}) \ill{} \ill{}
\note{l'arc ainsi défini est souligné par des traits dans le texte}

\struck{Il y a aussi des « arcs de spécialisation mutuellement disjoints » exceptionnels, qui \ill{}.}
\note{la fin de cette ligne biffée, « mutuellement disjoints », est reportée au-dessus, non biffée}

On va aussi définir des \emph{arcs} de spécialisation exceptionnels, \uncertain{qui} \uncertain{correspondent} aux couloirs de $\Sigma \cup D$ adjacents à $D$ comme \ill{} \ill{} et tels que le sommet du couloir soit sommet de $\Sigma$, \struck{$L$ et que} \add{et que} $L$ \struck{\ill{}} \add{\ill{} \ill{}} rencontre l'intérieur du \struck{\ill{}} \uncertain{couloir} \add{\ill{}} \struck{\ill{}} \ill{} \ill{} la rive. \struck{\ill{}} Le couloir la longe de $\partial\Delta = \widetilde{D}$, qui est un des deux arcs \uncertain{ouverts} de $L$, délimités par la classe d'équivalence de cordes définie par $s$. Ce n'est pas un arc de spécialisation ordinaire (car les \struck{\ill{}} cordes extrêmes sont dans la $=$ classe \struck{\ill{}} des $D_i$) et on voit \uncertain{même} que cet arc ne rencontre aucun arc de spécialisation \add{ordinaire} ni d'arc de générisation.
\note{les lettres $L$ et $D$ alternent dans cet alinéa, pour la même position relative. Figures : au milieu, une droite $L$ (vert clair et vert foncé) passant par un sommet $s$ où se coupent $D_1$ et $D_2$ (orange), traversée d'autres pseudo-droites, une petite flèche le long de $L$ et, dans une accolade, « la rive » ; au bas, un disque au bord vert (clair et foncé) traversé de cordes orange, cinq régions voisines du bord hachurées en jaune, les extrémités de $D_1$ et $D_2$ marquées en haut et en bas du bord par de petits segments bleus}

\page{36}
Il peut arriver que $L$ admette \emph{deux} arcs de spécialisation exceptionnels, ce qui est le cas ssi $L$ est \emph{\uncertain{position}} « bissectrice » d'un couloir de $\mathcal{C}$. Dans ce cas, il n'existe pas d'autre arc de spécialisation, et exactement deux arcs de générisation, et l'ensemble de ces quatre arcs est l'ensemble des \struck{\ill{}} \add{arêtes} d'une structure de polygone topologique sur $\widetilde{D}$.
\note{figure : un disque au bord vert traversé de cordes orange ; en haut et en bas du bord, deux petits segments bleus où aboutit un faisceau de cordes presque parallèles}
\marginal{arcs caractéristiques quand \uncertain{sont} les \uncertain{arêtes} d'un polygone \uncertain{caractéristique}}
\note{note écrite en biais dans la marge, reliée par un trait à la rubrique c) qui suit}

Dans tous les cas, \struck{\ill{}} deux \emph{arcs de transition} (i.e.\ de spécialisation ou de générisation) \add{distincts} \struck{\ill{} \ill{}} ne se rencontrent que si l'un d'eux est un arc de spécialisation exceptionnelle, l'autre un arc de générisation, et alors ils sont adjacents.

On appelle \emph{arcs d'inertie} \add{(ordinaires)} (les \struck{\ill{}} adhérences des composantes connexes du complémentaire de la réunion de l'ensemble des arcs de transition. En leur équivalent combinatoire en termes des sommets sur $\widetilde{L}$. Mais,
\note{« arcs d'inertie » est souligné ; la parenthèse n'est pas refermée}

\struck{Ainsi \ill{}} on appelle \emph{sommets de transition} les sommets sur $\widetilde{L}$ qui sont extrémités d'arcs de transition. Ces sommets partagent $\widetilde{L}$ en des arcs, qui sont de trois types :
\begin{enumerate}
\item[a)] les arcs de générisation
\item[b)] les arcs de spécialisation (ordinaires et exceptionnelles)
\item[c)] les \struck{arcs de} arcs qui ne sont ni de type a), ni de type b), et qu'on appelle \add{\uncertain{arcs} caractéristiques} (ils sont « \uncertain{inertes} » \uncertain{des} \emph{arcs d'inertie} les \uncertain{précurseurs} des transitions d'une pos.\ rel.\ à une autre)
\end{enumerate}
\note{« sommets de transition » : « de transition » est souligné ; les trois rubriques sont réunies par une accolade dans la marge}

Dans le cas où il n'y a pas d'arc de spécialisation exceptionnel, on pour deux arcs \add{caractéristiques} (adjacents), \struck{\ill{}} un est un seul est un arc d'inertie, l'autre étant \uncertain{dans} \uncertain{des} spécialisations. Deux arcs \uncertain{de} générisation, ou de spécialisation ne sont jamais adjacents. \uncertain{D'}inertie ne sont jamais adjacents, \uncertain{si} deux arcs caractéristiques sont adjacents, \uncertain{si} \uncertain{l'un} et \uncertain{si} \uncertain{l'aucun} \ill{} \struck{\ill{}} n'est un arc d'inertie, alors, \struck{\ill{}} \struck{\ill{} \ill{} \ill{}} l'un est de spécialisation exceptionnelle, l'autre est un arc de générisation. Parfois, le arc pour le $\mathcal{C}$ se réduit au sommet qui
\note{la fin de la colonne, rapide et raturée, est lue en partie ; « \struck{Deux} » biffé devant « arcs de générisation »}

\page{37}
commun, i.e.\ \struck{tout} \add{un} sommet qui est extrémité d'un arc de spéc.\ exceptionnel, s'appelle arc d'inertie exceptionnel (par abus de langage). Il peut être considéré de l'élargir, un \uncertain{peu} \uncertain{forçant} un peu la définition de l'arc de spécialisation exceptionnel, et le faisant s'arrêter aux premiers sommets \uncertain{suivant} ses extrémités primitives (disons ses \struck{\ill{}} constitutives « initiales » des \struck{\ill{}} arêtes extrêmes \add{\ill{}} de l'arc de spécialisation initial) : \struck{\ill{}} $\widetilde{L}$ \ill{} de l'arc de spécialisation (exceptionnel) « rétréci ».

Le bénéfice, c'est que les \add{transitions} \struck{spécialisations} i.e.\ les arcs de \struck{spécialisation} \add{transition} sont mutuellement disjoints dans tous les cas, et \add{de} deux arcs caractéristiques \uncertain{rétrécis} adjacents, \struck{\ill{}} un et un seul est arc de \uncertain{neutralité} (au sens « rétréci »), les arcs de faire d'exception.
\note{figure à gauche du texte : un disque au bord vert (clair et foncé) traversé de cordes orange, cinq régions voisines du bord hachurées en jaune, deux petits segments bleus en haut et en bas du bord}

Comme le nom l'indique, les arcs de générisation servent : \struck{\ill{}} \add{\uncertain{\ill{}}} une opération de générisation, qui consiste à « écarter les brins » au-dessus dudit arc, de façon que les deux sommets antipodiques cessant d'être arcs de générisation
\note{figures : trois disques au bord vert, traversés de cordes orange, reliés par des flèches $\mapsto$ et $\leftarrow$ ; sur le premier, deux cordes $D$, $D'$ se rejoignent en haut et en bas en deux petits segments verts ; le deuxième est couvert de hachures au crayon ; sur le troisième les extrémités de $D$ et $D'$ sont écartées en haut, croisées en bas ($D'$, $D$), un petit triangle orange au bas du bord}

(\emph{NB} Le polygone structural $\widetilde{L}$, à isom.\ comb.\ canonique près, n'est pas changé par \uncertain{cette} \uncertain{écriture} en question est \uncertain{duale} par arc de spécialisation. Dans $X$, l'opération se visualise comme un \emph{décollement} \add{\uncertain{\ill{}}} du sommet $s$ de $\Sigma$ dans la direction \uncertain{voulue} : celle où se trouve l'arc de générisation, qui « pousse » sur $L$, et en suivant $L$, donne l'arc de spécialisation correspondant.
\note{« décollement » est souligné ; « antibalayage », écrit au-dessus, est d'une lecture \uncertain{douteuse}. Figure : deux pseudo-droites orange $D$, $D'$ se croisant sur une droite verte (pleine et pointillée), un petit segment foncé au point de croisement, une flèche vers le bas ; $\rightleftarrows$ ; la même figure, le croisement passé de l'autre côté de la droite verte}

Ceci dit, si $(L, a)$ et $(L', a')$ \add{(a arc de gén.\ de $L$, $a'$ arc de spécialisation de $L'$)} se correspondent, il est utile d'\emph{attacher} $\Delta_L$ \emph{à} $\Delta_{L'}$, par identification de $a$ avec $a'$, \ill{} par un \ill{} \uncertain{les} \ill{}
\note{la parenthèse ouverte après « $(L, a)$ » n'est pas refermée. Figure à gauche : un disque au bord vert traversé par deux cordes orange qui se rejoignent au bas du bord, prolongé par des pointillés où elles se croisent au-dessous}

\page{38}
qui applique \add{\uncertain{l'}}extrémité de l'un des \struck{\ill{}} arcs, sur l'unique extrémité de l'autre qui appartienne à la « même » \struck{pseudo-droite} corde i.e.\ pseudo-droite, donc en \struck{\ill{}} \add{\ill{}} \uncertain{intervertissant} les extrémités, donnant \uncertain{lieu} à la \struck{figure} figure ci-contre.
\note{figure : deux disques au bord vert traversés de cordes orange, régions voisines du bord coloriées en jaune ; sur chacun, les extrémités $u$, $v$ en haut (une flèche au-dessus), et deux cordes $D$, $D'$ ; « $+$ » entre eux, puis « $=$ » et une accolade qui renvoie à un troisième disque accolé au premier par le bas, les extrémités $s$, $t$ et $t$, $s$ se touchant au point de contact, de petites flèches courbes marquant l'échange des extrémités, et en bas $v$, $u$}

En le faisant pour les deux arcs de générisation antipodiques, on trouve la figure ci-dessous. Mais réflexion faite, comme le recollement là correspond à un \uncertain{isom.}\ de transport entre $\underline{\omega}(\Delta_L) \simeq \underline{\omega}(L)$ et $\underline{\omega}(\Delta_{L'}) \simeq \underline{\omega}(L')$ qui est l'isomorphisme « évident », il me semble que ce n'est pas le bon \struck{\ill{}} « tranche » l'isom.\ d'identification qu'il faudrait plutôt \struck{\ill{}} prendre, \uncertain{mais} celui qui transforme $s$ en $s$, $t$ en $t$, de qui \uncertain{\ill{}} $\simeq$ \ill{} fait en \uncertain{intervertir} $D$ par $D'$, $D'$ par $D$ ; de façon à \uncertain{\ill{}} faire la symétrie par \uncertain{à} « l'axe » $(s, t) = a$,
\note{figures : un disque ($D'$, $D$ en haut) accolé par le bas, en un petit segment foncé où $D$ et $D'$ se croisent, à un second disque ($D$, $D'$), puis celui-ci à un troisième ($D'$, $D$ au bas) ; deux disques isolés à droite, reliés par des flèches au crayon aux points de contact}

\note{figure, en haut de la moitié droite : quatre disques au bord vert, traversés de cordes orange (pleines et pointillées), accolés deux à deux en quatre points de contact (petits segments foncés), autour d'un domaine central bordé d'arcs bruns ; les extrémités des cordes sont marquées $D_0$, $D'_0$, $D_1$, $D'_1$, et les coins voisins des contacts portent les lettres $\alpha$, $\beta$, $\beta'$}
Par cette identification toutes les « faces » sont d'ordre $4$, où les sommets de la \struck{\ill{}} carte déduite de celle déduite du graphe des positions relatives, avec la semi-circularisation définie par les polygones caractéristiques et les isom.\ de transport par l'orientation qui est l'opposé de l'isom.\ tautologique. Les coins obtenus ici comme faces, qui sont déduits des coins formés par les arcs d'inertie, ont des sommets, correspondant aux arêtes des « coins caractéristiques », qui sont de trois types, suivant les types des arcs de transition qui les bordent :
\note{la phrase « où les sommets \ldots{} tautologique » semble incomplète ; « carte déduite de celle déduite » est bien écrit ainsi}

\page{39}
\add{1 côté des} \struck{type} spécialisation (les arcs \struck{qui} \add{de transition} sont de spécialisation)

\add{1 côté des} \struck{type} générisation (les arcs de transition sont de gén.)

1 \struck{côté du type} mixte.
\note{« côté des » est écrit au-dessus de « type » biffé, aux deux premières lignes ; « côté du type » est souligné et biffé à la troisième}

Quand on contourne le coin en partant du côté spécialisation, on trouve dans l'ordre :

1) spécialisation $\to$ spécialisation

2) générisation $\to$ spécialisation

3) générisation --- générisation

4) spécialisation $\to$ générisation
\note{la deuxième rubrique porte un chiffre peu net, lu 2) ; à la troisième, un trait sans pointe}

(\struck{Opérer} \struck{\ill{}} dans un sens, ou dans l'autre) \struck{\ill{}} \uncertain{donc} il y a \struck{deux côtés} \add{une arête} du type spéc., une autre du type générisation, et deux du type mixte. Les faces correspondent biunivoquement aux \struck{\ill{}} \add{couples} $(L, a)$, \struck{où} $L$ est une position relative, et $a$ un arc \struck{de spécialisation} d'inertie du type spécialisation (ou du type générisation, au choix, il y a corr.\ 1-1 entre les deux types de couples). Le deuxième type est peut-être plus joli comme description, car il s'agit de trouver
\begin{enumerate}
\item[a)] les paires \struck{de sommets} $(s, t)$ \add{(critiques \uncertain{fermées})} \uncertain{\ill{}} \add{\uncertain{\ill{}}} de $\Sigma$ et les p.r.\ qui y \uncertain{mènent} ;
\item[b)] Pour \struck{\ill{}} telles données $(s, t, L)$,
\end{enumerate}
\note{la rubrique a) est corrigée à plusieurs reprises : « $(s, t)$ » est écrit au-dessus, « (critiques \uncertain{fermées}) » souligné, un mot \ill{} avant « de $\Sigma$ »}

\note{figure : un disque au bord vert, traversé de cordes orange, deux faisceaux de cordes se croisant ; autour du bord, quatre arcs fléchés numérotés $1$, $2$, $3$, $4$ au crayon}
s'il y a un arc de neutralité ayant \struck{\ill{}} ses extrémités au-dessus de $(s, t)$ (a priori, il est parmi les quatre arcs ci-contre\ldots) et les choisir. C'est suivant une description (très) particulière à expliciter, le cas particulier où les deux sommets \uncertain{désignés} forment une unique position relative.

La construction de la surface cellulaire n'est pas complète, car on a laissé de côté, dans le processus de recollement, les positions relatives supercritiques, et les arcs de spécialisation et de générisation qui vont aux ou partent de ces dernières. Donc \struck{\ill{}} il manque encore des disques polygonaux qu'il faut recoller sur les \uncertain{arêtes} qui restent libres pour ce recollement : on trouve \uncertain{ainsi}
\note{« arc de neutralité » : le mot, peu net, est lu comme à la page 37 --- lecture \uncertain{probable}}

\page{40}
\uncertain{ment}, par le procédé incomplet précédent en recollant le long des arcs de spécialisation exceptionnelle, une \emph{surface à bord} \add{\uncertain{cellulaire}} (compacte), chose qui se présente quand on a un graphe $\Gamma = (S, \vec{A}, \sigma, \sigma)$ pourvu des \struck{bords libres} \emph{arêtes libres} (i.e.\ \add{des} éléments \add{de $\vec{A}$} laissés fixes par $\sigma$), qui \struck{\ill{}} \add{\uncertain{peuvent}} s'interpréter comme des « arêtes orientées sans extrémité » (n'allant nulle part). Quand on forme les disques polygonaux $\operatorname{Dis}_s$, et qu'on les recolle par des isom.\ d'identification pour les $\bar{a}$ \emph{ordinaires} ($\sigma\bar{a} \neq \bar{a}$), on trouve toujours une surface à bord.
\note{« $\Gamma = (S, \vec{A}, \sigma, \sigma)$ » : les deux dernières lettres sont semblables sur la page ; la graphie $\bar{a}$ (barre simple) est la sienne ici}

Ici il s'agit de boucher les trous qui \struck{restent}, \struck{\ill{}} une fois bouchés les « trous quadrangulaires » qui viennent d'être (définis par des circuits de $4$ arcs d'inertie), par les disques correspondant aux positions relatives \emph{supercritiques}, qu'il s'agit de recoller \struck{\ill{}} suivant les arcs de spécialisation exceptionnels, correspondant ici aux arêtes libres. On est donc amené à étudier les \emph{circuits} formés d'\struck{arêtes} \add{arcs} de spécialisation exceptionnels, et voir comment les boucher.
\note{« qui viennent d'être » est suivi d'un participe peu net ; « supercritiques » est souligné}

\struck{Plus} Il sera plus simple peut être de s'inspirer de ce qui a déjà été fait, pour \uncertain{\ill{}} commençons l'étude au cas des spécialisations et pos.\ rel.\ exceptionnelles ($=$ supercritiques). Tout d'abord, pour une \struck{\ill{}} $L = D_i \in \Sigma$, on définit encore $\Delta_L$, en considérant $D_i$ comme p.r.\ \uncertain{de} p.r.\ : $\underline{\Sigma} = (X, \Sigma \smallsetminus \lbrace D_i \rbrace)$, ---
\note{le $\Sigma$ doublement souligné désigne ici le système privé de $D_i$, sans signe distinctif sur la page. La page s'arrête sur un tiret ; la suite est à la page 41 (lot 3)}

\end{document}
